% =====================================================================
%  Finding looping tracks in ACTS: a segment-and-merge architecture
%  Design note.  Line references are against ACTS 664c82a5 (19 Aug 2026).
% =====================================================================
\documentclass[11pt,a4paper]{article}

\usepackage[margin=2.4cm]{geometry}
\usepackage{newtxtext,newtxmath}
\usepackage{graphicx}
\usepackage{booktabs}
\usepackage{amsmath}
\usepackage{xcolor}
\usepackage[colorlinks=true,linkcolor=blue!45!black,urlcolor=blue!45!black,
            citecolor=blue!45!black]{hyperref}
\usepackage{enumitem}
\usepackage{caption}

% Works whether pdflatex is run from navigator_studies/ (via the Makefile)
% or from tex/ directly.
\graphicspath{{figures/}{../figures/}}

\captionsetup{font=small,labelfont=bf,width=0.92\textwidth}
\setlist[itemize]{itemsep=2pt,topsep=3pt}
\setlist[enumerate]{itemsep=2pt,topsep=3pt}

\newcommand{\code}[1]{\texttt{\small #1}}
\setlength{\emergencystretch}{2.5em}
\newcommand{\srcref}[2]{\code{#1:#2}}

\title{\bfseries Finding looping tracks in ACTS\\
       \large A segment-and-merge architecture}
\author{Design note}
\date{19 August 2026 \quad---\quad ACTS \code{664c82a5}}

\begin{document}
\maketitle

\begin{abstract}
\noindent
Low-$p_T$ tracks in a solenoidal field curve back towards the beamline instead
of leaving the detector. The stock ACTS \code{Navigator} cannot follow them
through the radial turning point, because the navigation layer extrapolates
along the tangent and is curvature-blind by interface. This note argues against
building a navigator that flies through the turn, and in favour of
\emph{cutting} the trajectory at the turning point and merging the resulting
segments afterwards. The central observation is that consecutive radial
turning points are separated by exactly half a revolution of the track circle,
so every segment is monotonic in $r$ and sub-half-helix in length --- precisely
the regime in which the existing navigator is already correct. The note
identifies the concrete extension points (\code{branchStopper} rather than a new
\code{navigator\_t}), proposes matching on helix invariants rather than on
bound parameters, and recommends an implementation order.
\end{abstract}

\tableofcontents

% =====================================================================
\section{The problem}
% =====================================================================

A charged particle in a uniform solenoidal field describes a helix. In the
transverse plane this projects to a circle of radius
\begin{equation}
  R \;=\; \frac{p_T}{0.3\,B}
  \qquad (\text{$R$ in\,m}, \; p_T \text{ in\,GeV}, \; B \text{ in\,T}),
\end{equation}
centred at some point $\vec{c}$. The track therefore never exceeds
\begin{equation}
  r_{\max} \;=\; |\vec{c}| + R ,
  \label{eq:rmax}
\end{equation}
which for a particle originating at the beamline ($|\vec{c}| = R$) is $2R$. If
$r_{\max}$ falls inside the tracking volume, the particle turns around and comes
back in. It is a \emph{looper}.

These tracks are not exotic. In a compact, high-field tracker a substantial
fraction of the soft charged multiplicity loops. They are, however, largely
invisible to combinatorial track finding as currently configured, and the reason
is architectural rather than physical.

\subsection{Why the stock navigator cannot follow the turn}

The \code{CombinatorialKalmanFilter} (CKF) is a propagator \emph{actor}. It
never talks to the \code{Navigator} directly. The two communicate only through
the stepping loop in \srcref{Propagator.ipp}{99--193}, and the entire contract
between them is the single field \code{state.navigation.currentSurface}
(Figure~\ref{fig:loop}).

\begin{figure}[htbp]
  \centering
  \includegraphics[width=0.95\textwidth]{fig-01-propagator-loop.pdf}
  \caption{The three-way handshake. The Navigator and the CKF never call each
    other; \code{Propagator.ipp} mediates. A null \code{currentSurface} is
    indistinguishable, to the CKF, from ``nothing here yet''.}
  \label{fig:loop}
\end{figure}

Three properties of this arrangement conspire against loopers.

\paragraph{The navigation layer is curvature-blind by interface.}
\code{Navigator::nextTarget(state, position, direction)} receives a
\emph{direction}, not a momentum. Gen3's \code{NavigationArguments} carries
\code{position}, \code{direction} and \code{tolerance} --- no $q/p$, no charge,
no field. Correspondingly, \code{Surface::intersect} is straight-line only:
there is no helix--surface intersection anywhere in ACTS Core. The navigator
therefore proposes targets by extending the tangent, and the error of that
estimate grows as (bend per step) $\times$ (step length), as sketched in
Figure~\ref{fig:chord}.

\begin{figure}[htbp]
  \centering
  \includegraphics[width=0.62\textwidth]{fig-02-tangent-vs-chord.pdf}
  \caption{Tangent, chord and true arc at increasing curvature per step. Panel
    (c) is the degenerate case: at the turning point the tangent is purely
    azimuthal and the radial gradient vanishes.}
  \label{fig:chord}
\end{figure}

\paragraph{At the turning point the extrapolation leaves the reachable region.}
Write $\hat{n} = (x,y)/r$ for the outward unit \emph{transverse} radial vector.
Where $\hat{d}\cdot\hat{n} = 0$, the tangent ray is purely azimuthal. Every
surface it meets at larger radius lies outside $r_{\max}$ and is
\emph{physically unreachable} by this track, while the surfaces the track will
actually cross --- the ones it already passed, at smaller radius --- lie behind
the ray and are invisible to a forward search. Figure~\ref{fig:turn} shows the
geometry.

\begin{figure}[htbp]
  \centering
  \includegraphics[width=0.86\textwidth]{fig-03-turning-point.pdf}
  \caption{The failure in the transverse plane. At $T$ the tangent scores
    candidates along a ray that leaves the reachable region
    $r \le r_{\max} = |\vec{c}| + R$; the layers the track will really
    re-cross are behind it.}
  \label{fig:turn}
\end{figure}

\paragraph{Gen3 never re-resolves inside a volume.}
Navigation candidates are resolved once per volume and consumed by a
monotonically increasing index (\srcref{Navigator.cpp}{574--586}). A refresh is
triggered only when \code{INavigationPolicy::isValid} returns \code{false}, and
the base implementation returns \code{true} unconditionally
(\srcref{INavigationPolicy.hpp}{265--271}) --- no concrete policy in Core
overrides it. So even if the track's direction reverses completely, the ordered
candidate list from volume entry is never rebuilt (Figure~\ref{fig:policy}).

\medskip\noindent
On top of this, loop protection caps the path length at
$\text{loopFraction} \times 2\pi p / B$, i.e.\ half a helix turn by default
(\srcref{PropagatorOptions.hpp}{46--48}), and \code{maxSteps = 1000}
(\srcref{PropagatorOptions.hpp}{34}) is a budget for the \emph{entire} CKF branch
tree, so a single runaway branch discards every track found so far.

% =====================================================================
\section{The pivot: cut, don't fly through}
% =====================================================================

The natural response to Section~1 is to build a curvature-aware
\code{navigator\_t}: a look-behind sweep along the chord, a curvature-bounded
layer window, a \code{checkTargetValid} that invalidates once the direction has
rotated. All of that is feasible --- see Appendix~\ref{app:navigator} --- but it
is a large amount of delicate new code in the most safety-critical part of the
propagation.

The alternative is to \textbf{stop at the turning point, and merge the pieces
afterwards}. This is strictly better, and the reason is geometric.

\subsection{Segments are half-revolutions}

Parametrise the transverse circle by the angle $\psi$ measured from the
direction of $\vec{c}$, so that a point on the track is
$\vec{x}(\psi) = \vec{c} + R\,(\cos\psi, \sin\psi)$ rotated into that frame.
Then
\begin{equation}
  r^2(\psi) \;=\; |\vec{c}|^2 + R^2 + 2\,|\vec{c}|\,R\cos\psi ,
  \label{eq:rpsi}
\end{equation}
and therefore
\begin{equation}
  \frac{\mathrm{d}(r^2)}{\mathrm{d}\psi} \;=\; -2\,|\vec{c}|\,R\,\sin\psi
  \;=\; 0
  \qquad\Longleftrightarrow\qquad
  \psi \in \{0,\ \pi\} .
\end{equation}
The stationary points of the radius are $\psi = 0$, giving
$r_{\max} = |\vec{c}| + R$, and $\psi = \pi$, giving
$r_{\min} = \bigl||\vec{c}| - R\bigr|$. \textbf{Consecutive radial turning
points are separated by $\Delta\psi = \pi$: exactly half a revolution of the
track circle.}

Two consequences follow immediately.

\begin{enumerate}
\item \textbf{Every segment is monotonic in $r$.} Between turning points,
  $\mathrm{d}r/\mathrm{d}s$ does not change sign. The trajectory is exactly the
  kind of outward- (or inward-) going arc that the stock navigator handles
  correctly. None of the three failure modes of Section~1.1 is reachable inside
  a segment.
\item \textbf{Every segment fits inside the default loop-protection budget.}
  A half revolution has 3D path length $s = \pi R / \sin\theta = \pi p / (0.3 B)$,
  which is precisely $\text{loopFraction} \times 2\pi p / B$ at the default
  $\text{loopFraction} = 0.5$. The half-turn cap stops being an antagonist and
  becomes a backstop that should rarely fire before the explicit turning-point
  stop does.
\end{enumerate}

\noindent
This is the whole argument for the architecture: \emph{by cutting at the turning
point, one never enters the regime that breaks the navigator.} The detection
criterion is also directly the geometry. With $\hat{n} = (x,y)/r$ the outward
unit transverse radial vector,
\begin{equation}
  \hat{d}\cdot\hat{n}
  \;=\; \frac{x\,d_x + y\,d_y}{r}
  \;=\; \frac{\mathrm{d}r}{\mathrm{d}s} ,
  \label{eq:radcos}
\end{equation}
the rate of change of the transverse radius per unit \emph{3D} path length.
Normalising the transverse direction instead, $\hat{d}_T = (d_x,d_y)/\sin\theta$,
gives the true cosine $\hat{d}_T\cdot\hat{n} = \mathrm{d}r/\mathrm{d}s_T \in
[-1,1]$. The two differ only by the positive factor $\sin\theta$, so they share
their sign and their zeros, and either may be used for detection.

From Eq.~\eqref{eq:rpsi}, $\mathrm{d}r/\mathrm{d}\psi \propto -\sin\psi$, so
Eq.~\eqref{eq:radcos} vanishes exactly at $\psi \in \{0, \pi\}$ --- at the
turning points and nowhere else.

% =====================================================================
\section{Components}
% =====================================================================

The proposed scheme has four parts. They are presented here in dependency
order, which is not the order in which they were first enumerated.

\subsection{Turning-point detection --- \code{branchStopper}}
\label{sec:stopper}

This is neither a navigator concern nor an aborter. The correct hook is the
\code{branchStopper} delegate declared in
\srcref{CombinatorialKalmanFilterExtensions.hpp}{60--61}. It returns
\begin{center}
\code{Continue} \quad\textbar\quad \code{StopAndDrop} \quad\textbar\quad
\code{StopAndKeep}
\end{center}
(\srcref{CombinatorialKalmanFilterExtensions.hpp}{38--42}). It is invoked on both
paths through the actor: the measurement path at
\srcref{CombinatorialKalmanFilter.hpp}{870} and the hole/material path at
\srcref{CombinatorialKalmanFilter.hpp}{733}. Returning \code{StopAndKeep} stores
the branch, pops it from \code{activeBranches}, and lets the depth-first tree
walk continue normally (Figure~\ref{fig:tree}). No propagation surgery is
required, and no new \code{navigator\_t} is needed.

\begin{figure}[htbp]
  \centering
  \includegraphics[width=0.92\textwidth]{fig-05-ckf-branch-tree.pdf}
  \caption{The CKF branch walk. \code{StopAndKeep} stores and pops; it does
    \emph{not} restart from the turning point, which is why the scheme is
    necessarily multi-pass.}
  \label{fig:tree}
\end{figure}

Four implementation notes.

\paragraph{Use the transverse radial cosine.} Compute $\hat{d}\cdot\hat{n}$ of
Eq.~\eqref{eq:radcos} --- the projection onto the \emph{transverse} radial
$\hat{n} = (x,y)/r$ --- and \textbf{not} the projection onto the spherical
radial $\hat{r} = \vec{x}/|\vec{x}|$. The latter,
\begin{equation*}
  \hat{d}\cdot\hat{r} \;=\; \frac{x\,d_x + y\,d_y + z\,d_z}{|\vec{x}|} ,
\end{equation*}
carries the extra term $z\,d_z$, which for a forward-going track is strictly
positive. It can therefore stay positive straight through the transverse
turning point, and the flip is never seen.

Only the \emph{sign} is needed, and $r > 0$, so the division may be dropped
entirely:
\begin{center}
\code{const double radialProj = pos.x() * dir.x() + pos.y() * dir.y();}
\end{center}
Store it per branch (below) and stop when it changes sign.

\paragraph{Guard the small-$r$ case.} $\hat{n}$ is ill-defined as $r \to 0$, so
the sign of \code{radialProj} becomes noise-dominated for a track passing close
to the beamline --- which is exactly the $r_{\min}$ turning point of a primary
looper, where $r_{\min} = \bigl||\vec c| - R\bigr| \approx 0$. Require $r$ above
a few mm before trusting a flip, or handle inner turning points separately.

\paragraph{Carry the previous value in a dynamic column.} Sign-flip detection
needs one step of history, and that history must be \emph{per branch}: two
branches off the same surface can turn at different points. The pattern is
already in the tree --- \srcref{TrackFindingAlgorithm.cpp}{176--255} attaches a
\code{BranchState} struct via \code{ProxyAccessor} on a named column. This
survives branch splits: \code{copyFromShallow} calls
\code{copyFromWithoutStates}, which calls \code{copyDynamicFrom} at
\srcref{TrackProxy.hpp}{551}.

\paragraph{Do not let the existing stopper eat short segments.}
\srcref{TrackFindingAlgorithm.cpp}{247} returns \code{StopAndDrop} whenever
\code{!enoughMeasurements}. A track already curving back as it enters the barrel
yields a two- or three-hit segment which is legitimately below
\code{minMeasurements} but is still a real piece of a real track.
\textbf{Segment-level retention must be looser than track-level retention}, with
the tight quality cut moved to the merged object. This is plumbing, but it is
the kind that silently deletes exactly the signal one is chasing.

\subsection{Continuation: directed extrapolation vs.\ independent seeding}
\label{sec:continuation}

Because \code{StopAndKeep} stores and pops rather than restarting, segment~2 is
a \emph{separate} \code{findTracks} invocation. There are two ways to obtain its
starting state, and the choice reshapes the rest of the design.

\paragraph{(a) Independent seeding, then matching.} Seed every segment on its
own, find them all, then pair them up. This is the more general scheme, and it
is the only way to recover a segment whose predecessor was never found. It is
also expensive: an inbound segment does not point at the beamline, so
\emph{triplet seeding rejects it on impact parameter long before $p_T$ matters}.
The blocker is \code{impactMax}, not \code{minPt}, and loosening \code{impactMax}
enough to accept inbound segments is a combinatorial disaster.

\paragraph{(b) Directed continuation.} At the turning point the fitted
parameters and their full covariance are already known, and so is the helix.
Extrapolate analytically across the turn --- a single closed-form step, no
search, no navigator --- and hand the result to the next CKF pass as its
starting state. The continuation is \emph{constructed} rather than discovered.

\medskip\noindent
Directed continuation should be the primary mechanism. It removes the need for
inbound-capable seeding entirely for all segments after the first, and it makes
matching free for the pairs it produces, because the pairing is known by
construction. Independent seeding plus matching is then a \emph{fallback} for
orphan segments, not the main path.

The cost of (b) is error inheritance: a mismeasured first segment gives a
displaced continuation seed, and at low $p_T$ the extrapolation uncertainty
across the turn is large. That is tolerable --- one opens the search window wide
at the first layer of segment~2, which is exactly what the CKF is good at.

\subsection{Segment matching}
\label{sec:matching}

For the orphan case one still needs a matching criterion. The obvious
formulation --- Euclidean distance in bound-parameter space --- does not work:
\code{loc0} and \code{loc1} are in mm, \code{phi} and \code{theta} in radians,
\code{qOverP} in $\mathrm{GeV}^{-1}$. There is no defensible relative weighting,
and the covariances differ by orders of magnitude between a soft and a stiff
segment. Two better metrics, used at different stages:

\paragraph{Pre-selection on helix invariants.} Both segments lie on the
\emph{same circle}. Compute per segment
\begin{equation}
  \bigl(c_x,\ c_y,\ R,\ \tan\lambda,\ q\bigr)
\end{equation}
and match on those. They are physically commensurate, directly comparable, and
--- crucially --- require no propagation to a common surface, so pre-selection
can be done by binning in $O(N \log N)$ rather than by $O(N^2)$ propagations.
The shared circle also supplies a common parametrisation $\psi$ for free, so the
longitudinal match reduces to comparing $z$ extrapolated to a common $\psi$.

There is a bonus. Energy loss makes $R \propto p_T$ shrink monotonically along
the track, so $R_2 < R_1$ is not noise to be tolerated --- \textbf{the sign of
the radius mismatch identifies which segment came first}. That halves the
pairing combinatorics and provides an independent consistency check.

\paragraph{Confirmation by combined refit.} Once a pair survives pre-selection,
do not try to make the distance test precise. Refit the union of the hits with
the Kalman filter and let the combined $\chi^2/\mathrm{ndf}$ arbitrate. Matching
then only has to be good enough to \emph{propose}, which is a far weaker
requirement than being good enough to \emph{decide}. If a distance test is
wanted in addition, use the Mahalanobis form
\begin{equation}
  \chi^2 \;=\; \Delta p^{\mathsf T}\,\bigl(C_1 + C_2\bigr)^{-1}\,\Delta p ,
\end{equation}
which is $\chi^2$-distributed on 5 degrees of freedom and therefore gives a
principled threshold, rather than an arbitrary distance cut.

\paragraph{Trap.} Segment~2 travels inward. If it was seeded independently, its
direction and charge conventions may be reversed relative to segment~1 --- a
sign-flipped $q/p$ and a $\phi$ offset by $\pi$ on genuine matches. Normalise
conventions before comparing anything.

\medskip\noindent
There is no segment-merging machinery in Core to build on. The classes under
\code{Acts/\allowbreak AmbiguityResolution/} solve duplicate \emph{removal}, not
concatenation. The merger is therefore new code --- but it lives entirely
downstream of the propagation, operating on a finished
\code{TrackContainer}, which makes it easy to develop and test standalone.

\subsection{Widening acceptance}
\label{sec:widening}

This should come \emph{last}, once the preceding parts make the effect
measurable.

For the measurement selector: \code{MeasurementSelectorCuts}
(\srcref{MeasurementSelector.hpp}{34--42}) bins on \code{etaBins} only. There is
no $p_T$ binning, so ``widen for loopers'' currently means ``widen for
everyone'', which costs branches and burns the shared \code{maxSteps} budget.

Before widening anything, \textbf{check the $\chi^2$ pulls} on a sample of known
low-$p_T$ tracks:
\begin{itemize}
\item If the pulls are wider than unit width, the predicted covariance is wrong
  --- the material model or the field integration is underestimating multiple
  scattering. Widening the cut then papers over a bug that will resurface in the
  fit.
\item If the pulls are unit-width and hits are still being lost, widen --- but as
  an explicit $p_T$- or $R$-dependent selector, not globally.
\end{itemize}

For seeding, see Section~\ref{sec:continuation}: under directed continuation the
seeding requirement for segments after the first disappears, which is the
largest single simplification available here.

\subsection{Disc-region navigation}
\label{sec:discs}

A dedicated \code{LoopNavigator} for the endcap region was proposed. It is worth
being precise about which problem it would solve, because in one reading it is
unnecessary:

\begin{itemize}
\item \emph{A track turns in the barrel and its continuation lies in the endcap
  discs.} This is just a normal monotonic segment that happens to start in a
  different volume. The stock navigator handles it. No new navigator is needed
  --- and this is the design working as intended.
\item \emph{A track spirals within the disc region}, crossing the same disc
  twice at small $r$. This is a genuine multi-intersection problem. But it has
  the same turning-point structure in a different coordinate, so the same
  segment-and-merge treatment applies with the cut placed on the sign flip of
  $\mathrm{d}z/\mathrm{d}s$ instead of $\mathrm{d}r/\mathrm{d}s$.
\end{itemize}

\noindent
Either way, the entire value of this architecture is that it keeps the
propagation inside the regime the existing navigator is correct in. Introducing
a custom navigator reintroduces the risk that was just designed away. It should
be held in reserve.

% =====================================================================
\section{Recommended order}
% =====================================================================

\begin{enumerate}[label=\textbf{\arabic*.}]
\item \textbf{Turning-point \code{branchStopper}} (Section~\ref{sec:stopper}).
  Roughly sixty lines. It requires no changes to the propagator, the navigator
  or the geometry, and it immediately yields a measurable population of turning
  points to characterise. Do this first, and instrument it.
\item \textbf{Directed continuation} (Section~\ref{sec:continuation}). Analytic
  extrapolation across the turn, feeding a second CKF pass. This is what makes
  the scheme work without touching seeding.
\item \textbf{Segment matching for orphans} (Section~\ref{sec:matching}). Helix
  invariants for pre-selection, combined refit for confirmation.
\item \textbf{Acceptance widening} (Section~\ref{sec:widening}), with pull
  distributions in hand rather than by guesswork.
\item \textbf{A custom navigator} (Section~\ref{sec:discs}, Appendix~%
  \ref{app:navigator}) --- only for whatever residual remains after measuring
  steps 1--4. If that residual is small, the hardest part of the project has
  been avoided entirely.
\end{enumerate}

% =====================================================================
\appendix
\section{If a custom navigator does turn out to be needed}
\label{app:navigator}
% =====================================================================

The design conclusions reached before the pivot, retained in case step~5 is
reached:

\begin{itemize}
\item \textbf{Separate ``no target this step'' from ``navigation has failed''.}
  This --- not its brute-force candidate set --- is the actual source of
  \code{TryAllNavigator}'s robustness. It returns \code{None} \emph{without}
  setting \code{navigationBreak} at \srcref{TryAllNavigator.hpp}{413}, letting the
  propagator step blindly forward and recover.
\item \textbf{Steal the chord-based look-behind sweep}
  (\srcref{TryAllNavigator.hpp}{340--346}, Figure~\ref{fig:lookbehind}). Search
  backwards along the chord with $\text{nearLimit} = -|AB| + \varepsilon$ and
  $\text{farLimit} = 0$. Roughly 35 lines. Drop the flat whole-volume candidate
  set that surrounds it in \code{TryAllNavigator}; it is Gen1-only and
  expensive.
\item \textbf{Measure curvature from two consecutive chords} rather than from
  $q/p$. This needs no interface change --- sidestepping the curvature-blindness
  of \code{NavigationArguments} --- and it is immune to the seed-stage $p_T$
  uncertainty that makes an analytic $r_{\max}$ cut (Eq.~\ref{eq:rmax})
  unusable in the first few layers.
\item \textbf{Replace ``current layer only'' with a curvature-bounded layer
  window}, and make the validity checks curvature-aware rather than constant:
  both \code{checkTargetValid} and the policy's \code{isValid}. Note that
  \code{isValid} governs \emph{cadence}, not
  \emph{scope}, so a policy override works regardless of how many volumes the
  loop crosses.
\item \textbf{Shrink the step size as the radial gradient approaches zero}, so
  the turn is not overshot and a block of layers skipped.
\end{itemize}

\begin{figure}[htbp]
  \centering
  \includegraphics[width=0.85\textwidth]{fig-04-chord-lookbehind.pdf}
  \caption{The retrospective sweep. The chord is the secant of the arc actually
    flown, known only \emph{after} the step --- aim with the tangent, then verify
    along the chord.}
  \label{fig:lookbehind}
\end{figure}

\begin{figure}[htbp]
  \centering
  \includegraphics[width=0.95\textwidth]{fig-06-policy-lifecycle.pdf}
  \caption{\code{INavigationPolicy} call sites across a volume traversal, and
    why \code{isValid} defaulting to \code{true} means Gen3 never re-resolves
    inside a volume.}
  \label{fig:policy}
\end{figure}

\clearpage
% =====================================================================
\section{Code reference table}
% =====================================================================

All line numbers against ACTS \code{664c82a5} (19 August 2026).

\begin{table}[htbp]
\centering
\small
\begin{tabular}{@{}llp{6.9cm}@{}}
\toprule
\textbf{File} & \textbf{Lines} & \textbf{What} \\
\midrule
\multicolumn{3}{@{}l}{\itshape Propagation} \\
\code{Propagator.ipp}   & 99--193 & the stepping loop \\
                        & 51      & \code{nextTarget} \\
                        & 56      & pre-step \code{updateSurfaceStatus} \\
                        & 102     & \code{step} \\
                        & 129     & post-step \code{updateSurfaceStatus} \\
                        & 135     & \code{handleSurfaceReached} \\
                        & 144     & \code{actorList.act} \\
                        & 166     & \code{checkTargetValid} \\
\code{PropagatorOptions.hpp} & 34 & \code{maxSteps = 1000} (whole branch tree) \\
                        & 40      & \code{maxTargetSkipping = 100} \\
                        & 46--48  & \code{loopProtection}, \code{loopFraction = 0.5} \\
\midrule
\multicolumn{3}{@{}l}{\itshape Navigation} \\
\code{Navigator.cpp}    & 335, 560 & \code{isValid} call sites \\
                        & 380     & \code{popState} on volume exit \\
                        & 574--586 & candidate index; no mid-volume re-resolve \\
                        & 645     & \code{createState} on volume entry \\
                        & 672     & \code{initializeNavigationCandidates} \\
\code{INavigationPolicy.hpp} & 265--271 & \code{isValid} defaults to \code{true} \\
\code{TryAllNavigator.hpp} & 340--346 & chord look-behind sweep \\
                        & 413     & \code{None} without \code{navigationBreak} \\
\midrule
\multicolumn{3}{@{}l}{\itshape Track finding} \\
\code{CombinatorialKalmanFilterExtensions.hpp} & 38--42 & \code{BranchStopperResult} enum \\
                        & 60--61  & \code{BranchStopper} delegate signature \\
                        & 67      & default (always \code{Continue}) \\
\code{CombinatorialKalmanFilter.hpp} & 471--538 & \code{Actor::reset} \\
                        & 733     & \code{branchStopper}, hole/material path \\
                        & 870     & \code{branchStopper}, measurement path \\
\code{MeasurementSelector.hpp} & 34--42 & \code{MeasurementSelectorCuts}; \code{etaBins} only \\
\code{TrackProxy.hpp}   & 551     & \code{copyDynamicFrom} --- columns survive splits \\
\code{TrackFindingAlgorithm.cpp} & 176--255 & example \code{BranchStopper} \\
                        & 247     & \code{StopAndDrop} on \code{!enoughMeasurements} \\
\bottomrule
\end{tabular}
\end{table}

\paragraph{Facts worth restating.}
\begin{itemize}
\item \code{Surface::intersect} is straight-line only. There is no
  helix--surface intersection anywhere in ACTS Core.
\item \code{NavigationArguments} carries \code{position}, \code{direction} and
  \code{tolerance} --- no $q/p$, charge or field.
\item \code{maxSteps} is a budget for the entire branch-tree walk; exhausting
  it returns \code{Propagation\allowbreak Reaches\allowbreak MaxSteps} and discards every track found
  so far.
\item Dynamic track columns are copied on branch split, so per-branch state is
  safe.
\end{itemize}

\end{document}
